\ssscp{Parallelisms and doublets}{02-03-02-03}
Throughout the region in all of our Wallacean subgroups,
both parallelisms and doublets are found to varying degrees,
not only in poetic ritual language,
but also in everyday speech
(\citealp{Fox1974,Fox1988,Fox2014,Fox2016}, \citealp[15--31]{GrimesTherikGrimesJacob1997}, \citealp[162--164]{TaberTaber2015}).
Parallelisms tend to be clausal,
while doublets (or word pairs, as illustrated in the next paragraph)
can be hyphenated or joined in a
unitary phrase by a conjunction.
And the two can work in concert,
as illustrated by the translation of part of a \ili{Molo} \it{onen} ‘prayer’ in
\citet{GrimesBIP}: \it{“Lord,
when we encounter difficulty
and suffering, abuse and grumbling, you make us cherish love and inner peace.
Lord, when we encounter abuse
and grumbling, you make us cherish forgiveness.”}
\citet[379--385]{Edwards2020}
discusses the sophisticated interplay of metathesis involved with \ili{Amarasi} parallelisms.
The use of synonyms in parallelisms
and doublets throughout the region are part of what makes language
beautiful, and speakers who use them eloquent.

Whether looking at parallelisms or doublets,
most often similar categories of words or near synonyms are treated
as \it{synonymous} (same meaning) when used together.
Antonyms or \it{counterparts} often reflect whole categories:
mother-father = parents; night-day = around the clock; ascend-descend
= go all over; plate-fork = eating utensils; pigs-chickens = domesticated animals.
And sometimes the pair is \it{elaborative},
fleshing out the picture slightly,
but that is the rarest type.
It is quite common for verbs that have slightly different meaning
in isolation to be used synonymously,
such as \ili{Kotos} \ili{Amarasi} \it{t-kumain ma t-humaʔmoe} ‘we laugh
and we smile’, \it{n-poi ma n-boor} ‘it appears and bursts forth’.
Sometimes these synonymous pairs are easiest to gloss simply
as ‘cut1-cut2’, or ‘carry1-carry2’.
Fox (various) has written extensively about additional figurative
or symbolic meanings when used in ritual language registers (\it{bini})
in the languages of Rote.
An example from nearby Timor,
in the \ili{Kotos} \ili{Amarasi} \it{aʔa sramat} genre,
the phrase \it{mutiʔ ma mnatuʔ} (lit: ‘silver
and gold’) is conventionally used to refer to ‘kings
and nobles’, whereas \it{koro-manu} (lit: bird-chicken) is conventionally
used to refer to ‘populace, commoners’ \citep{GrimesBIP}.
\citet[28]{GrimesTherikGrimesJacob1997} observe “In some pairs,
the meaning of one unit may not be recoverable.
There are several instances where the only occurrence of a word
in the entire available data corpus is in doublet form,
and native speakers \it{cannot} give a meaning for the word in
isolation.” These are more commonly the second member of a doublet
and are often identifiable as being from neighbouring languages or dialects.
\citet{Valeri1990} calls this “the retreat of the semantics”
in central Maluku, in favour of celebrating links
and alliances with other languages
and ethnic groups. \citet[30]{GrimesTherikGrimesJacob1997}
summarise, “The skilful use of parallelisms is at once
a link to the ancestors,
a link expressing solidarity with those of other dialects or
languages with whom one is allied,
and a social link with those who hear
and appreciate the skilful art of the poet.” These patterns of
using doublets and parallelisms to celebrate synonyms,
and taking synonyms from other dialects or languages are well
documented for many languages
and subgroups in the region.
Since the use of parallelisms is also documented in the Philippines
and Borneo, there is no reason to assume it is a recent development
restricted to this region.
\citet{GrimesBIP} also describes similar patterns found in Papuan
languages, such as the TAP language \ili{Klon} on \ili{Alor}.

